\section{数据供应链：从数据源到能力的工业流程}

如果把数据看成产业，就需要一条供应链：采集、授权、清洗、标注、合成、混合、训练、评估、反馈。每一环都会影响最终模型能力。数据供应链越长，越需要标准化、审计和质量控制。

\subsection{数据供应链表}

\noindent\begin{longtable}{p{0.20\textwidth}p{0.34\textwidth}p{0.37\textwidth}}
\toprule
环节 & 任务 & 质量风险 \\
\midrule
采集 & 从网页、机器人、仿真、用户行为获取数据 & 来源偏差、版权、隐私。 \\
清洗 & 去重、去噪、过滤低质样本 & 误删高价值长尾。 \\
标注 & 给任务、动作、状态、反馈加结构 & 标注不一致、成本高。 \\
合成 & 用模型或仿真补充稀缺场景 & 分布偏移、模式崩塌。 \\
混合 & 确定数据配比和训练顺序 & 能力偏科、遗忘。 \\
评估 & 检查模型能力和失败类型 & 指标失真、过拟合 benchmark。 \\
回流 & 把失败和新场景带回数据池 & 闭环慢、反馈不准。 \\
\bottomrule
\end{longtable}

\begin{knowledgebox}{数据供应链的壁垒}
数据壁垒不只在“我有数据”，还在“我知道怎样把数据变成能力”。清洗、配比、评估和回流，常常比原始数据更难复制。
\end{knowledgebox}

\subsection{本章小结}

数据产业的工业化，意味着从数据源到能力之间有一整条供应链。任何一环薄弱，都会降低最终模型能力。

